\chapter{Research Synthesis}
\label{sec:synthesis}

Transformer-based \ac{VLA} models have ushered robotics into an era of true quasi-generalist capability, demonstrating that a single model can, to an extent, see, reason and act across diverse tasks; yet leading examples \cite{PI0.52025, HELIX2025, GR00T2025, GEMER2025}, despite enabling zero-shot manipulation of novel objects, remain far from the master-of-all-trades that real-world deployment demands. By weaving together the threads of classical control, \ac{AI}, transformer-based foundation models, \ac{RL} and modular architectures, we synthesize these advances to shape our research approach, leveraging breadth from pre-trained models, depth from online adaptation and structure from modular design, and we highlight the gaps that must be bridged for truly adaptive robotic learning in \ac{OOD} environments.

\section{Key Insights}

\paragraph{Foundation Models Enable Breadth.}
Foundation models provide robots with a priori knowledge and generalization capabilities previously unseen in robotics. Models like RT-2 and PaLM-E demonstrate that a single model can leverage web-scale data to perform many tasks and even exhibit zero-shot behaviors \cite{RT22023,PALM2023}. This breakthrough allows handling \emph{novel inputs} (new objects or instructions) because the model likely encountered similar examples during training (e.g., RT-2 recognizing a ``coffee mug'' it never saw in fine-tuning). However, foundation models alone are not a panacea: they tend to be heavy, inflexible after training, and their real-world performance can degrade under distribution shift without further adaptation \cite{COLOSSEUM2024}.

\paragraph{RL Enables Depth.}
\ac{RL} gives robots the ability to improve with experience in new scenarios. \ac{RL} can fine-tune a policy to dynamics or reward specifics not covered in pre-training \cite{RLSURVEY2013}. It is essential for \emph{online learning}: when encountering an \ac{OOD} environment, \ac{RL} can adjust the policy incrementally. Yet \ac{RL} alone is too slow and unsafe to tackle large distribution shifts from scratch. The synergy is clear: a foundation model supplies a strong initial policy, and \ac{RL} refines that policy to the specific context efficiently \cite{PPO2017}, achieving \emph{fast adaptation} with orders of magnitude fewer trials than learning from zero.

\paragraph{Hybrid and Modular Architectures Provide a Structural Prior.}
By not relying on a single monolithic neural network, hybrid and modular architectures enforce safety, handle multi-objective trade-offs, and incorporate specialist knowledge. For example, a classical \ac{PID} loop can maintain low-level stability, and safety checks can override learned policies if a collision is imminent. Structured models often outperform end-to-end approaches in robustness \cite{CHRL2019}. Moreover, modular design aligns with practical engineering: individual modules (e.g., vision systems) can be upgraded without retraining the entire policy, which is vital for long-term system development.

\paragraph{Research Hypothesis.}
An approach that integrates a transformer-based foundation model (for generalization), \ac{RL} (for adaptation), and a modular architecture (for safety and structure) is a promising path to adaptive robotic learning. Alone, each component has limitations: foundation models without adaptation may fail under distribution shift, \ac{RL} without a strong prior is too slow, and rigid modular systems lack novelty handling. Together, they form a more complete solution.

From these insights, our research hypothesis is reinforced: an approach that integrates a transformer-based foundation model (for generalization), \ac{RL} (for adaptation), and a modular architecture (for safety and structure) is a promising path to adaptive robotic learning. None of these components alone would likely succeed: foundation models without adaptation might fail when their assumptions break, \ac{RL} without a strong prior is too slow, and a rigid modular system without learning can't handle novelty. Together, they form a more complete solution.

\section{Outstanding Gaps and Challenges}

\paragraph{Generalization vs. Specialization Balance.}
Large foundation models excel at broad generalization but can lack the precision or efficiency of smaller, specialized policies trained via \ac{RL} for specific robots and tasks. To bridge this gap, we must develop methods such as task-specific fine-tuning of large models or distilling their knowledge into compact high-frequency controllers that preserve broad capabilities while delivering task-level performance. For example, one could distill RT-2’s policy into a lightweight network optimized for real-time execution. Our research will investigate fine-tuning and distillation strategies that adapt foundation models to particular robots and environments without sacrificing their general knowledge.

\paragraph{Safety and Trustworthiness.}
Modular architectures allow embedding safety rules, yet formally verifying that a learning-enabled robot remains safe under all real-world variations is generally intractable. Control-theoretic tools (e.g., Lyapunov stability analysis) offer partial guarantees, but neural policies defy exhaustive analysis. Practical approaches include safety shields and extensive simulation under perturbations-practices we will adopt. A critical open problem is online monitoring: detecting when the robot enters an unknown regime and gracefully reverting to a provably safe baseline. We will explore anomaly-detection modules (e.g., flagging out-of-distribution sensor inputs via a generative model) to trigger safe mode, acknowledging this remains a challenging research frontier.

\paragraph{Continual Learning and Memory.}
While our current framework enables fast adaptation in a new environment, it does not ensure retention of prior skills. As a robot encounters ever more varied settings, it should accumulate competencies without catastrophic forgetting. True continual learning is unsolved: naively fine-tuning via \ac{RL} will overwrite earlier knowledge. Potential solutions include modular compartmentalization of new skills (e.g., adding task-specific subnetworks) or meta-learning schemes that consolidate parameters. We identify the stability–plasticity dilemma: maintaining performance on environment A while learning B, as a key gap. We plan experiments to measure forgetting and to test methods such as Elastic Weight Consolidation or orthogonal gradient descent for mitigating it.

\paragraph{Benchmarking and Evaluation.}
There is currently no standard benchmark for ``adaptive robotics.'' In this thesis we take a first step by introducing a structured real-robot evaluation protocol, but more broadly we argue that the community should develop a comprehensive evaluation suite drawing on modified Meta-World tasks, Colosseum-style real-world perturbation tests, and sensor-failure scenarios to measure adaptation across sensor, environment, and task changes. Such a protocol, spanning new objects, varied terrain, and shifting goals, would provide breadth and consistency and could serve as a shared benchmark for fair comparison.

\section{Synthesis for Approach}

Considering all of the above, we synthesize the approach as follows: We will start with a strong foundation (a pretrained transformer model endowed with multimodal understanding) and anchor it to a reliable control backbone (ensuring baseline competence and safety). On this foundation, we will layer an adaptation mechanism (\ac{RL} fine-tuning or online learning) that refines performance in novel scenarios. The expected outcome is a robot that, upon encountering an \ac{OOD} situation, initially behaves reasonably (not randomly, thanks to the foundation model) and then improves rapidly through experience (thanks to \ac{RL}), all while never venturing into catastrophic behavior (thanks to safety modules and constraints).

By iteratively testing this system on increasingly challenging distribution shifts, we will validate the core hypothesis. Success will be measured by the robot’s ability to maintain high task success with minimal drop after a shift and how quickly it can recover any lost performance. For example, if introducing a 30\% perturbation in the weight of an object causes the success to drop from 90\% to 50\%, can our system learn to return, say, 80\% + success in a small number of trials? Does it avoid dangerous attempts during that adjustment?

The literature gives optimism that each piece can achieve parts of these goals: foundation models have shown the ability to do things like pick up unfamiliar objects correctly on the first try, and meta-\ac{RL} has shown fast improvement within tens of trials on new tasks in simulation. The novel contribution of our work will be to seamlessly integrate these pieces in real robotic hardware and demonstrate end-to-end that the robot can take what it ``knows'' (from pre-training) and adapt what it does not (through experience), in a unified system.

In the next section, we transition from what needs to be done to how we will do it. We outline the concrete research questions we will answer and the experimental methodologies, metrics, and implementation steps that will operationalize this synthesized approach.